% ============================================================
%  Chapter 2 -- Regular Languages and Decision Problems
% ============================================================

\chapter{Regular Languages and Decision Problems}
\label{ch:regular-decision}

\aside{Regular languages are the finite-word
behaviours recognized by finite automata.
The point of the theory is not only
recognition, but effective verification.}

In \cref{ch:fsa} we built finite automata to recognize finite-word
behaviours: a turnstile event sequence, a parity check over $\{a,b\}$,
and so on.  Automata are an \emph{operational} description---they tell
us how to scan a word.  But they are not the only way to describe the
same languages, and they are not always the most convenient one.

This chapter adds three complementary viewpoints, each answering a
different question that the operational model alone leaves open.
\begin{itemize}[nosep]
  \item \emph{Can we describe a language algebraically, without drawing
  a graph?}  Regular expressions (\cref{sec:regular-expressions}) give
  such a syntax, and Kleene's theorem shows they define exactly the
  regular languages.
  \item \emph{If we are given a property built from simpler ones by
  union, intersection, complement, can we still build an automaton for
  it?}  Closure properties (\cref{sec:regular-closure}) answer yes, and
  the constructions are the backbone of every verification algorithm
  in the rest of the book.
  \item \emph{Once automata encode properties, can we decide inclusion,
  equivalence, universality?}  The decision problems of
  \cref{sec:finite-decision-problems} all reduce to a single
  reachability check---a finite-word prototype of model checking.
\end{itemize}
The chapter closes with the Myhill--Nerode theorem
(\cref{sec:myhill-nerode}), which explains \emph{why} some languages
admit a finite-state description at all and yields the canonical minimal
DFA.  The guiding theme throughout is \emph{reduction}: we reduce the
hard-looking semantic questions to reachability-style emptiness checks
that we actually know how to perform.


% ============================================================
\section{Regular Expressions}
\label{sec:regular-expressions}
% ============================================================

Drawing a DFA is an operational act: we specify how the machine moves.
For many purposes---communicating a property, composing specifications,
searching text---a more algebraic syntax is preferable, one that builds
the language out of pieces using a few combinators.  This is what
regular expressions provide.

\begin{definition}[Restricted regular expressions]
\label{def:restricted-regex}
The \textbf{restricted regular expressions} over an alphabet $\Sigma$
are generated by:
\[
  \emptyset,\qquad \varepsilon,\qquad a\in\Sigma,
\]
and, if $r$ and $s$ are regular expressions, by:
\[
  r+s,\qquad rs,\qquad r^*.
\]
\end{definition}

The semantics is the expected one.  The expression $r+s$ denotes union,
$rs$ denotes concatenation, and $r^*$ denotes finite repetition.

\begin{example}[A simple regular expression]
\label{ex:regex-even-a}
Over $\Sigma=\{a,b\}$, the expression
\[
  b^*(ab^*ab^*)^*
\]
denotes the language of words with an even number of occurrences of
$a$.  The initial $b^*$ consumes any number of $b$'s before the first
$a$, and each repetition of $ab^*ab^*$ contributes two occurrences of
$a$.
\end{example}

This is the same language recognized by the DFA in
\cref{fig:even-a-dfa}.  The expression describes the structure of the
word; the automaton describes how to scan it.

\begin{formaldetails}[Restricted, general, and star-free expressions]
The word \emph{restricted} is used because the syntax above contains
only union, concatenation, and Kleene star.  A more liberal syntax may
also include Boolean operations such as intersection and complement.
This does not add expressive power: once we can translate expressions to
automata, closure under intersection and complement lets us compile the
Boolean operations away and return to a restricted expression.

There is, however, an important smaller fragment.  A \textbf{star-free}
expression may use union, concatenation, complement, and intersection,
but not Kleene star.  Star-free languages are regular, but not every
regular language is star-free.  Later, when logic enters the course,
this fragment becomes visible as the finite-word shadow of first-order
definability.
\end{formaldetails}

Before proving the equivalence between expressions and automata, we
need a more flexible automaton model that lets us wire pieces together
without consuming input.  The wiring device is the
\emph{$\varepsilon$-move}: a transition that changes state while the
reading head stays put.  Without it, splicing two automata for
concatenation would force us to merge states by hand and keep track of
spurious symbols; with it, we simply add free edges and let a later
clean-up pass dissolve them.

\begin{definition}[$\varepsilon$-NFA]
\label{def:epsilon-nfa}
An \textbf{$\varepsilon$-NFA} is an NFA whose transition function may
also use the empty word $\varepsilon$ as a label.  Formally, the
transition function has the shape
\[
  \Delta : Q\times(\Sigma\cup\{\varepsilon\})\to 2^Q,
\]
so a transition in $\Delta(q,\varepsilon)$ changes state without
consuming any input symbol.
\end{definition}

Operationally, an $\varepsilon$-NFA behaves like an NFA with one extra
freedom: at any moment it may follow an $\varepsilon$-labelled edge and
jump to another state, while the input pointer does not move.  This is
purely a construction device---it gives us the wiring freedom needed to
compose automata for union, concatenation, and Kleene star---but, as we
show next, it does not add any expressive power.

To short-circuit every chain of free moves, we need to know, for each
state, where it can reach by following only $\varepsilon$-edges.  This
is the $\varepsilon$-closure.

\begin{definition}[$\varepsilon$-closure]
\label{def:epsilon-closure}
For a state $q$, the \textbf{$\varepsilon$-closure}
$\mathsf{ECl}(q)$ is the set of states reachable from $q$ by taking
zero or more $\varepsilon$-moves.  For a set of states $S$,
\[
  \mathsf{ECl}(S)=\bigcup_{q\in S}\mathsf{ECl}(q).
\]
\end{definition}

Because $Q$ is finite, $\mathsf{ECl}(q)$ is the least fixpoint of the
monotone operator $X\mapsto\{q\}\cup\bigcup_{p\in X}\Delta(p,\varepsilon)$.
This is the same fixpoint flavour that will return, scaled up, in the
symbolic reachability algorithms of \cref{subsec:reachability-fixpoint}.

\begin{theorem}[$\varepsilon$-removal]
\label{thm:epsilon-removal}
For every $\varepsilon$-NFA $A$ there exists an NFA $A'$ without
$\varepsilon$-moves such that $L(A)=L(A')$.
\end{theorem}

\begin{proof}
Let $A=(Q,\Sigma,\Delta,q_0,F)$.  We define an ordinary NFA
$A'=(Q,\Sigma,\Delta',q_0,F')$ on the same states and alphabet.  The
construction encodes the slogan \emph{take all free moves, read one
symbol, take all free moves again} directly into the transition
function:
\[
  \Delta'(q,a)
  =
  \mathsf{ECl}\!\left(
    \bigcup_{p\in \mathsf{ECl}(q)} \Delta(p,a)
  \right).
\]
This is the best possible one-step simulation of $A$: starting from
$q$, the run in $A$ may first wander through $\varepsilon$-moves,
arrive at some state $p$, consume $a$, and then again wander through
$\varepsilon$-moves; $\Delta'(q,a)$ collects exactly the destinations
reachable that way.

There is one boundary case.  In $A$, the empty word is accepted if and
only if some accepting state is reachable from $q_0$ by $\varepsilon$-moves
alone, since the input pointer never moves.  The new automaton would
otherwise miss this case, because every transition of $A'$ consumes a
real symbol.  We correct the final set:
\[
  F'=
  \begin{cases}
    F\cup\{q_0\}, & \mathsf{ECl}(q_0)\cap F\neq\emptyset,\\
    F, & \text{otherwise.}
  \end{cases}
\]
A straightforward induction on the length of the input word $w$ now
shows that the set of states reached by $A'$ on $w$ from $q_0$ equals
the $\varepsilon$-closure of the set of states reached by $A$ on $w$.
Intersecting with $F'$ then yields exactly the same accepting verdict as
intersecting with $F$ in $A$.
\end{proof}

With $\varepsilon$-NFAs as a flexible intermediate model, we can equate
the syntactic and operational views of regularity.

\begin{theorem}[Kleene equivalence]
\label{thm:kleene}
A language $L\subseteq\Sigma^*$ is recognized by a finite automaton if
and only if it is denoted by a regular expression.
\end{theorem}

\begin{proof}
We prove both directions constructively.

\medskip
\noindent\textbf{Regular expression to automaton.}
We build an $\varepsilon$-NFA by structural induction on the expression.
For $\emptyset$, use a single non-accepting state with no transitions.
For $\varepsilon$, use an initial state that is already accepting.  For
a symbol $a$, use one transition labelled $a$ from the initial state to
a fresh accepting state.

Assume $\varepsilon$-NFAs have already been built for $r$ and $s$, with
disjoint state sets and a single accepting state each (this is without
loss of generality, since we can always redirect outgoing edges through
a fresh accepting sink via $\varepsilon$-moves).
\begin{itemize}[nosep]
  \item For $r+s$, create a fresh initial state with $\varepsilon$-moves
  to the initial states of the two automata, and $\varepsilon$-moves
  from each of their accepting states to a fresh accepting state.
  \item For $rs$, connect the accepting state of the automaton for $r$
  to the initial state of the automaton for $s$ by an $\varepsilon$-move;
  the accepting state of $s$ becomes the accepting state of the combined
  automaton.  \Cref{fig:regex-concat} shows the wiring.
  \item For $r^*$, add a fresh initial state that is already accepting,
  add an $\varepsilon$-move to the initial state of the automaton for
  $r$, and add an $\varepsilon$-move from the accepting state of $r$
  back to that initial state.
\end{itemize}
Each construction mirrors exactly the semantic operation: choose a
branch for union, run one machine after the other for concatenation, and
repeat any finite number of times for star.  By
\cref{thm:epsilon-removal} we can then eliminate $\varepsilon$-moves,
determinize, and obtain a DFA.  Hence every regular expression denotes
an automaton-recognizable language.

\medskip
\noindent\textbf{Automaton to regular expression.}
Let $A$ be a DFA with states $q_1,\ldots,q_n$, with $q_1$ initial.  For
$0\leq k\leq n$, define $R_{i,j}^k$ to be the language of all words that
take the automaton from $q_i$ to $q_j$ using only intermediate states
among $q_1,\ldots,q_k$.

We prove by induction on $k$ that each $R_{i,j}^k$ is denoted by a
regular expression.  For $k=0$, no intermediate state is allowed, so a
path from $q_i$ to $q_j$ is either the empty path (when $i=j$) or a
single transition labelled by some $a\in\Sigma$.  Hence $R_{i,j}^0$ is a
finite union of symbols, together with $\varepsilon$ when $i=j$; this is
regular.

For the induction step, a path counted by $R_{i,j}^k$ either never
visits $q_k$ as an intermediate state---in which case it lies in
$R_{i,j}^{k-1}$---or it visits $q_k$ at least once.  In the latter case,
split the path at the first visit to $q_k$, then at each complete return
to $q_k$, and finally at the last departure toward $q_j$.  This yields
the recurrence
\[
  R_{i,j}^k =
  R_{i,j}^{k-1}
  +
  R_{i,k}^{k-1}\,(R_{k,k}^{k-1})^*\,R_{k,j}^{k-1},
\]
which is regular by the induction hypothesis.

Finally, the language of the automaton is
\[
  L(A)=\bigcup_{q_j\in F} R_{1,j}^n,
\]
again a regular expression.
\end{proof}

\begin{figure}[ht]
\centering
\begin{tikzpicture}[
  node distance=1.9cm,
  state/.style={circle, draw, thick, minimum size=8mm, fill=blue!8},
  acc/.style={circle, draw, double, thick, minimum size=8mm, fill=teal!12},
  arr/.style={-Stealth, thick},
]
  \node[state] (ir) {$i_r$};
  \node[acc, right=of ir] (fr) {$f_r$};
  \node[state, right=of fr] (is) {$i_s$};
  \node[acc, right=of is] (fs) {$f_s$};
  \draw[arr] (ir) -- node[above, font=\small] {$r$} (fr);
  \draw[arr, orange!70!black, thick] (fr) -- node[above, font=\small] {$\varepsilon$} (is);
  \draw[arr] (is) -- node[above, font=\small] {$s$} (fs);
  \node[below=8mm of fr, font=\small, align=center]
    {The accepting exit of the first automaton is spliced to the\\
     entry of the second by a single $\varepsilon$-move (in orange).};
\end{tikzpicture}
\caption{The $\varepsilon$-NFA construction for concatenation in the
regex-to-automaton direction of \cref{thm:kleene}: the accepting state
of $r$ and the initial state of $s$ are wired together by an
$\varepsilon$-move, and the fresh accepting state $f_s$ inherits
acceptance.  The other two operators (union, Kleene star) follow the
same splicing idea.}
\label{fig:regex-concat}
\end{figure}

A useful consequence of \cref{thm:kleene}, combined with the closure
constructions we are about to develop in \cref{sec:regular-closure}, is
that the two syntactic classes introduced above---restricted and
general---recognize exactly the same languages.

\begin{corollary}[Restricted and general expressions coincide]
\label{cor:regex-classes}
Restricted and general regular expressions denote the same class of
languages.
\end{corollary}

\begin{proof}
Every restricted expression is, by definition, a general one.  Conversely,
given a general expression, we translate each Boolean operator into an
automaton construction: complementation becomes final-state swapping on
a complete DFA, and intersection becomes the product construction of
\cref{def:product-automaton}.  This produces a finite automaton, which
\cref{thm:kleene} converts back into a restricted expression.
\end{proof}

The star-free fragment is more subtle.  Removing Kleene star does cost
expressive power: there are regular languages that no star-free
expression can denote.  Our running example
\cref{ex:regex-even-a}---words with an even number of $a$'s---is the
canonical witness.  Its minimal DFA cycles between two states and is
not \emph{aperiodic}; by Sch\"utzenberger's theorem, which we will meet
in \cref{ch:deterministic-automata-star-free}, this is exactly the
obstruction that keeps it out of the star-free fragment.  Star-free
languages are therefore a \emph{strictly weaker} subclass of the regular
languages, and they will reappear later as the finite-word shadow of
first-order definability.


% ============================================================
\section{Closure Properties}
\label{sec:regular-closure}
% ============================================================

Closure properties say that regular languages remain regular after
standard operations.  They are useful because verification constantly
builds compound properties from simpler ones.

\begin{theorem}[Closure of regular languages]
\label{thm:regular-closure}
Regular languages over finite words are closed under union,
concatenation, Kleene star, complementation, and intersection.
\end{theorem}

Union, concatenation, and star follow directly from regular expressions
or $\varepsilon$-NFA constructions.  Complementation and intersection
are more important for verification, so we spell them out.

\subsection{Complementation}
\label{subsec:finite-complementation}

If $A=(Q,\Sigma,\delta,q_0,F)$ is a complete DFA, its complement is
\[
  \overline{A}=(Q,\Sigma,\delta,q_0,Q\setminus F).
\]
The same run is used; only the final verdict changes.

\begin{observation}
\label{obs:complete-deterministic-complement}
The simple final-state swap works only for complete deterministic
automata.  Non-determinism and missing transitions both break the
argument.
\end{observation}

For an NFA, one accepting branch is enough to accept a word.  Swapping
final and non-final states would not turn ``there exists an accepting
branch'' into ``all branches reject''.  The standard route is therefore:
determinize, complete, then swap final states.

\subsection{Intersection}
\label{subsec:finite-intersection}

Intersection is handled by running two automata in lockstep.

\begin{definition}[Product automaton]
\label{def:product-automaton}
Let
$A_1=(Q_1,\Sigma,\delta_1,q_1,F_1)$ and
$A_2=(Q_2,\Sigma,\delta_2,q_2,F_2)$ be DFAs.  Their product for
intersection is the DFA with:
\[
  Q=Q_1\times Q_2,\qquad
  q_0=(q_1,q_2),
\]
\[
  \delta((p_1,p_2),a)=(\delta_1(p_1,a),\delta_2(p_2,a)),
  \qquad
  F=F_1\times F_2.
\]
\end{definition}

The first component remembers where $A_1$ is; the second remembers where
$A_2$ is.  A word is accepted exactly when both components end in
accepting states.

The word \emph{end} is doing real work here.  For finite words, both
automata consume exactly the same input positions and there is one final
moment at which acceptance is tested.  This is why the simple accepting
set $F_1\times F_2$ is correct.  In \cref{ch:omega-automata}, the same
idea must be modified: two B\"uchi automata may visit their accepting
states at different times, so simultaneous final-state membership is no
longer enough.

\begin{figure}[ht]
\centering
\begin{tikzpicture}[
  cell/.style={draw, minimum width=18mm, minimum height=10mm, align=center},
  good/.style={fill=teal!12},
  bad/.style={fill=gray!8},
]
  \matrix (m) [matrix of nodes, nodes={cell}, row sep=-\pgflinewidth, column sep=-\pgflinewidth] {
    $p_0,r_0$ & $p_0,r_1$ & $p_0,r_2$ \\
    $p_1,r_0$ & $p_1,r_1$ & $p_1,r_2$ \\
  };
  \node[draw, rounded corners=3pt, fit=(m-2-2), inner sep=2pt, fill=teal!12] {};
  \node[below=6mm of m, font=\small] {A product state stores one state from each automaton.};
\end{tikzpicture}
\caption{The state space of a product automaton is a Cartesian product.}
\label{fig:product-grid}
\end{figure}


% ============================================================
\section{Decision Problems}
\label{sec:finite-decision-problems}
% ============================================================

Once languages are represented by automata, many semantic questions
become graph questions.

\begin{definition}[Core decision problems]
\label{def:finite-decision-problems}
For finite automata over $\Sigma^*$, the standard decision problems are:
\begin{itemize}[nosep]
  \item \textbf{membership}: given $A$ and $w$, is $w\in L(A)$?
  \item \textbf{emptiness}: is $L(A)=\emptyset$?
  \item \textbf{universality}: is $L(A)=\Sigma^*$?
  \item \textbf{inclusion}: is $L(A_1)\subseteq L(A_2)$?
  \item \textbf{equivalence}: is $L(A_1)=L(A_2)$?
\end{itemize}
\end{definition}

Membership is simulation.  Emptiness is reachability: does some final
state occur on a path from the initial state?

\begin{proposition}[Emptiness for finite automata]
\label{prop:finite-emptiness}
Let $A$ be an NFA or DFA.  Then $L(A)\neq\emptyset$ if and only if some
accepting state is reachable from $q_0$ in the underlying directed graph.
\end{proposition}

\begin{proof}
If an accepting state is reachable, the labels along that path form an
accepted word.  Conversely, an accepted word induces a run from $q_0$ to
an accepting state.  Thus non-emptiness and accepting-state reachability
are the same question.
\end{proof}

There is another way to say the same thing, useful later when we move to
symbolic model checking: if a finite automaton accepts something, then it
accepts a short witness.

\begin{theorem}[Small model property for finite automata]
\label{thm:small-model-finite}
Let $A$ be a DFA with $n$ states.  Then
\[
  L(A)\neq\emptyset
  \quad\Longleftrightarrow\quad
  \text{there exists }w\in L(A)\text{ with }|w|<n.
\]
\end{theorem}

\begin{proof}
The right-to-left implication is immediate: an accepted word of any
length proves non-emptiness.

For the other direction, assume $L(A)\neq\emptyset$ and choose an
accepted word $z=a_1a_2\cdots a_m$ of minimum length.  Let
\[
  p_0 \xrightarrow{a_1} p_1
      \xrightarrow{a_2} \cdots
      \xrightarrow{a_m} p_m
\]
be the accepting run, with $p_0=q_0$ and $p_m\in F$.

If $m\geq n$, then the sequence $p_0,p_1,\ldots,p_m$ contains at least
$n+1$ state occurrences but only $n$ possible states.  By the pigeonhole
principle, $p_i=p_j$ for some $0\leq i<j\leq m$.  Split the word as
\[
  z = u v r,
\]
where $u=a_1\cdots a_i$, $v=a_{i+1}\cdots a_j$, and
$r=a_{j+1}\cdots a_m$.  The segment $v$ labels a loop from $p_i$ back to
$p_i$.  Removing that loop gives the shorter word $ur$, and the run
still goes from $q_0$ to the same accepting state $p_m$.  Thus $ur$ is
accepted, contradicting the minimality of $z$.  Hence the minimal
accepted word has length $m<n$.
\end{proof}

\begin{figure}[ht]
\centering
\begin{tikzpicture}[
  state/.style={circle, draw, thick, minimum size=9mm, fill=blue!8},
  acc/.style={circle, draw, double, thick, minimum size=9mm, fill=teal!12},
  arr/.style={-Stealth, thick},
]
  \node[state] (p0) {$p_0$};
  \node[state, right=1.7cm of p0] (pi) {$p_i$};
  \node[state, right=2.2cm of pi] (pj) {$p_j$};
  \node[acc, right=1.9cm of pj] (pm) {$p_m$};
  \draw[arr] (p0) -- node[above] {$u$} (pi);
  \draw[arr, bend left=28, red!70!black, dashed] (pi) to node[above] {$v$} (pj);
  \draw[arr] (pj) -- node[above] {$r$} (pm);
  \draw[blue!65!black, dashed, thick] (pi) to[bend right=40] node[below] {$p_i=p_j$} (pj);
\end{tikzpicture}
\caption{If an accepting run repeats a state, the loop can be removed to obtain a shorter accepted word.}
\label{fig:small-model-loop-removal}
\end{figure}

The graph algorithm is better in practice, but the small-model theorem is
conceptually useful: decidability follows because, in principle, we only
need to inspect finitely many words shorter than $n$.

\subsection{Reachability as a Fixpoint}
\label{subsec:reachability-fixpoint}

Let $\mathsf{Post}(X)$ be the set of states reachable in one transition
from some state in $X$.  Starting from $R_0=\{q_0\}$, define
\[
  R_{i+1}=R_i\cup \mathsf{Post}(R_i).
\]
Because $Q$ is finite, this increasing sequence eventually stabilizes at
a set $R^*$.  The language is non-empty exactly when
\[
  R^*\cap F\neq\emptyset.
\]
This fixpoint view is the finite-state ancestor of the symbolic
preimage/fixpoint algorithms used later for CTL and planning.

The other problems reduce to emptiness by using closure.

\begin{formaldetails}[Reductions]
\[
  L(A)=\Sigma^*
  \quad\Longleftrightarrow\quad
  L(\overline{A})=\emptyset.
\]
\[
  L(A_1)\subseteq L(A_2)
  \quad\Longleftrightarrow\quad
  L(A_1)\cap \overline{L(A_2)}=\emptyset.
\]
\[
  L(A_1)=L(A_2)
  \quad\Longleftrightarrow\quad
  L(A_1)\subseteq L(A_2)\text{ and }L(A_2)\subseteq L(A_1).
\]
\end{formaldetails}

For universality, we build the complement and ask whether any rejected
word exists.  For inclusion, we search for a counterexample:
\[
  w\in L(A_1)\setminus L(A_2).
\]
If the intersection $L(A_1)\cap\overline{L(A_2)}$ is empty, then no such
counterexample exists.  Equivalence is just two inclusions, because two
sets are equal precisely when neither contains an element missing from
the other.

This is the finite-word prototype of model checking.  The system
automaton describes possible behaviours; the specification automaton
describes allowed behaviours.  Verification asks for language inclusion.


% ============================================================
\section{The Myhill--Nerode View}
\label{sec:myhill-nerode}
% ============================================================

Automata theory also tells us when a language has a finite-state
representation at all.  The key idea is indistinguishability by
continuations.

\begin{definition}[Right congruence of a language]
\label{def:myhill-nerode-relation}
For a language $L\subseteq\Sigma^*$, define $x\equiv_L y$ if and only
if, for every continuation $z\in\Sigma^*$,
\[
  xz\in L \quad\Longleftrightarrow\quad yz\in L.
\]
\end{definition}

Two prefixes are equivalent when no future suffix can separate them.
After reading either prefix, the automaton needs exactly the same
information to decide all possible continuations.

\begin{definition}[Right-invariant equivalence]
\label{def:right-invariant-equivalence}
An equivalence relation $\sim$ on $\Sigma^*$ is \textbf{right-invariant}
if
\[
  x\sim y \quad\Longrightarrow\quad xz\sim yz
  \quad\text{for every }z\in\Sigma^*.
\]
\end{definition}

The relation $\equiv_L$ is right-invariant by definition: if no
continuation separates $x$ and $y$, then no longer continuation can
separate $xz$ and $yz$ either.

\begin{example}[The automaton relation]
\label{ex:automaton-right-congruence}
Let $A=(Q,\Sigma,\delta,q_0,F)$ be a DFA.  Define
\[
  x\equiv_A y
  \quad\Longleftrightarrow\quad
  \widehat{\delta}(q_0,x)=\widehat{\delta}(q_0,y).
\]
This relation has at most $|Q|$ classes.  It is right-invariant because
once two prefixes put the automaton in the same state, appending the
same suffix makes the two future computations identical.
\end{example}

\begin{theorem}[Myhill--Nerode theorem]
\label{thm:myhill-nerode}
A language $L\subseteq\Sigma^*$ is regular if and only if
$\equiv_L$ has finite index, meaning it partitions $\Sigma^*$ into
finitely many equivalence classes.
\end{theorem}

\begin{proof}
Assume first that $L$ is recognized by a DFA
$A=(Q,\Sigma,\delta,q_0,F)$.  If two words reach the same state, then
every continuation has the same fate from that state.  Formally, if
$\widehat{\delta}(q_0,x)=\widehat{\delta}(q_0,y)$, then for every
$z\in\Sigma^*$,
\[
  \widehat{\delta}(q_0,xz)=\widehat{\delta}(q_0,yz).
\]
Therefore $xz\in L$ if and only if $yz\in L$, so $x\equiv_L y$.  The
automaton relation $\equiv_A$ has at most $|Q|$ classes, and each
$\equiv_A$-class sits inside one $\equiv_L$-class.  Hence $\equiv_L$ has
finite index.

Conversely, assume $\equiv_L$ has finite index.  Build the quotient DFA
\[
  A_L=(\Sigma^*/{\equiv_L},\Sigma,\delta,[\varepsilon],F_L)
\]
where
\[
  \delta([w],a)=[wa]
  \quad\text{and}\quad
  F_L=\{[w]:w\in L\}.
\]
The transition is well-defined because $\equiv_L$ is right-invariant:
if $w\equiv_L w'$, then $wa\equiv_L w'a$.  The accepting set is
well-defined because taking the empty continuation $z=\varepsilon$ in
the definition of $\equiv_L$ shows that equivalent words are either both
in $L$ or both outside $L$.

After reading a word $w$, the quotient automaton is in state $[w]$.
Thus it accepts exactly the words whose equivalence class belongs to
$F_L$, which is exactly $L$.
\end{proof}

\begin{proposition}[Minimality]
\label{prop:myhill-nerode-minimality}
The quotient automaton built from $\equiv_L$ has the minimum possible
number of states among all DFAs recognizing $L$.
\end{proposition}

\begin{proof}
Let $A$ be any DFA recognizing $L$.  If two words reach the same state
in $A$, then they are $\equiv_L$-equivalent by the first half of the
previous proof.  Therefore a single state of $A$ can represent only one
Myhill--Nerode class.  Since every class is reached by some word, $A$
must have at least as many states as there are $\equiv_L$-classes.  The
quotient automaton has exactly that many states, so it is minimal.
\end{proof}

The minimal DFA for $L$ is therefore canonical up to renaming of states:
merge exactly the prefixes with the same continuation behaviour, and no
more.

\begin{summarybox}
\begin{itemize}[nosep]
  \item Regular expressions and finite automata define the same finite-word languages.
  \item Complementation needs complete deterministic automata.
  \item Intersection is implemented by a product construction.
  \item Inclusion and equivalence reduce to emptiness.
  \item Myhill--Nerode explains why regular languages are exactly those with finitely many distinguishable prefixes.
\end{itemize}
\end{summarybox}

\section*{Exercises}
\addcontentsline{toc}{section}{Exercises}

\begin{chapex}[Complementation needs completeness]
Draw a partial DFA over $\{a,b\}$ that accepts only the word $a$.  If you
swap final and non-final states without adding a sink state, what
language do you get?  Compare it with the true complement of $\{a\}$.
\end{chapex}

\begin{chapex}[Product construction]
Let $A_1$ recognize words ending in $a$, and let $A_2$ recognize words
with an even number of $a$'s.  Describe the accepting states of the
product automaton for $L(A_1)\cap L(A_2)$.
\end{chapex}

\begin{chapex}[Myhill--Nerode classes]
For the language of words over $\{a,b\}$ ending in $a$, list the
different continuation behaviours of prefixes.  How many states should
the minimal DFA have?
\end{chapex}
